\documentclass[t,12pt]{beamer}

\usepackage{amssymb}
\usepackage{amsfonts}
\usepackage{amsmath}
\usepackage{dsfont}
\usepackage[utf8x]{inputenc}
\usepackage[labelfont=bf,justification=justified]{caption}
\usepackage{subcaption}
\usepackage{fixltx2e}
\usepackage{setspace}
\usepackage{arydshln}
\usepackage{xcolor}
\usepackage{tikz}
\usepackage{chngpage}
\usepackage{comment}
\usetikzlibrary{shapes,arrows}
\tikzstyle{block} = [rectangle, draw, fill=blue!20, 
    text width=5em, text centered, rounded corners, minimum height=4em]
\tikzstyle{line} = [draw, -latex']

\usetheme{Emory}
\makeatletter
\setbeamertemplate{footline}
{
  \leavevmode%
  \hbox{%
  \begin{beamercolorbox}[wd=.333333\paperwidth,ht=2.25ex,dp=1ex,center]{author in head/foot}%
    \usebeamerfont{author in head/foot}\insertshortauthor~~\beamer@ifempty{\insertshortinstitute}{}{\insertshortinstitute}
  \end{beamercolorbox}%
  \begin{beamercolorbox}[wd=.333333\paperwidth,ht=2.25ex,dp=1ex,center]{title in head/foot}%
    \usebeamerfont{title in head/foot}\insertshorttitle
  \end{beamercolorbox}%
  \begin{beamercolorbox}[wd=.333333\paperwidth,ht=2.25ex,dp=1ex,right]{date in head/foot}%
    \usebeamerfont{date in head/foot}\insertshortdate{}\hspace*{2em}
    \insertframenumber{} / \inserttotalframenumber\hspace*{2ex} 
  \end{beamercolorbox}}%
  \vskip0pt%
}
\makeatother

\setbeamerfont{frametitle}{size=\large}
\setbeamertemplate{navigation symbols}{}

\titlegraphic{\includegraphics[width=0.3\textwidth]{~/Dropbox/Emory.png}}

\begin{document}

\title[Moving-average smoothing]{Moving-average smoothing}
\author{\vspace{0.35in} \\ {\fontfamily{ppl}\selectfont William Mann} \vspace{-0.35in}}

\date{}
\maketitle

	\begin{frame}[c]
	\vfill
	\centerline{Moving-average smoothing}
	\vfill
	\end{frame}

	\begin{frame}[c]{Features of economic series: Trend, cycle, residual}
	\begin{itemize}
	\item
	It is common to think of economic series as consisting of a long-term trend, a seasonal cycle, and a residual.
	\medskip
	\item
	A popular way to visualize and understand a data series is to try to decompose it into these components.
	\begin{itemize}
	\smallskip
	\item
	To be clear, there are many potential ways to do this. 
	\item
	They should all be interpreted carefully, as visualization tools, 
	\\ not necessarily as revealing any reliable statistical patterns.
	\end{itemize}
	\medskip
	\item
	The most common decomposition is based on moving averages.
	\\ We will overview this approach today.
	\end{itemize}
	\end{frame}

	\begin{frame}[c]{Centered moving averages}
	A \textbf{centered moving average} of order $m$ is as follows:
	\\ \medskip
	If $m$ is odd, then let $k = \frac{m-1}{2}$, and
	$$
	CMA_t \quad = \quad \frac{1}{m} \times (y_{t-k} + ... + y_{t+k})
	$$
	If $m$ is even, then let $k = \frac{m}{2}$, and
	$$
	CMA_t \quad = \quad \frac{1}{m} \times \left( \frac{1}{2} y_{t-k} + y_{t-k+1} + ... + y_{t+k-1} + \frac{1}{2} y_{t+k} \right)
	$$
	Centered moving averages are a common way to visualize the ``trend'' component of a data series.
	A common choice of $m$ is the order of seasonality in the data, for example $m=12$ for monthly data.
	\end{frame}

	\begin{frame}[c]{Seasonal component}
	After calculating the series of CMAs, a common approach to modeling the seasonal component of a time series is as follows:
	\begin{itemize}
	\medskip
	\item Calculate a ``detrended'' series as $y_{t,\text{detrend}} = y_t - CMA_t$.
	\smallskip
	\item For each stage $l$ within a cycle, calculate a seasonal average $\bar{y}_l$ as the average of $y_{t,\text{detrend}}$ across all the $l$ observations.
	\smallskip
	\item For example, if the data is monthly, then $\bar{y}_l$ is the average of $y_{t,\text{detrend}}$ for a specific month across all years in the sample.
	\smallskip
	\item The seasonal component for each stage $l$ is its seasonal average, minus the average of all the seasonal averages.
	\end{itemize}
	\medskip
	Finally, the residual is $y_t$ minus the trend and seasonal components.
	\\ \smallskip See example on Canvas.
	\end{frame}

	\begin{frame}[c]{Variations on the moving-average smoothing}
	The simple approach described above is very common in practice, 	
	\\ but you will also see many variations on it. Two examples are:
	\begin{itemize}
	\medskip
	\item \underline{Trailing averages}: When we want to build \emph{forecasts}, we don't have the future values and must average only over past values.
	\\ This often leads us to also use weighted averages (next point).
	\medskip
	\item \underline{Weighted averages}: In a two-sided moving average, it is common to weight the central observation by the most.
	In a trailing average, it is common to weight the \emph{most recent} observation by the most, to put the emphasis on the newest data.
	\end{itemize}
	\medskip
	Next week we will cover exponential smoothing, a popular approach that uses a trailing average with a specific weighting scheme.
	\end{frame}

\end{document}
